\section{Jobs: evidence and the pending access audit}\label{app:jobs}
\paragraph{Design and information.}
The primary simulation samples four occupations, 20 seed blocks and three
frozen months (August--October 2025), followed by two replication months
(November--December). It uses ACS 2020--2024 person records from the evaluation
donor partition and Revelio postings for the New York metropolitan area.
Each primary market contains 24 seekers and 24 postings, stratified by wage
group, with one opening per posting. Feasibility includes the specified pay
floor and seniority rule. Preferences and employer priorities are modeled;
the frozen wage weight is 2 and the seeker-noise scale is 0.75. Subsequent
calibration did not resolve the mismatch between ACS annual earnings and
posted salaries. These are conditional simulation results, not observed
applications, hiring, or effects of using an employment service.

The model receives structured occupation, experience proxy, education and pay
floor information, plus anonymized posting IDs, salary ranges and seniority.
It does not receive employer names, posting text or the seeker's wage group.
The recorded calls use \texttt{gpt-6-luna} through the approved HUIT route,
with reasoning disabled and at most 300 completion tokens. A dated model
snapshot is unavailable; the run manifest records returned model strings and
send times. Cached replies fix the evaluated menus.

\paragraph{Endpoints and paired comparisons.}
Let $M$ be full-list DA, $N_D$ the assignment from menu policy $D$, and
$B_M=\{i:M(i)\ne\bot\}$. The jobs coverage rate is
$|B_M|^{-1}\sum_{i\in B_M}\mathbf1\{(i,M(i))\in D\}$.
Placement uses all seekers: $|B|^{-1}\sum_i\mathbf1\{N_D(i)\ne\bot\}$.
The former is conditional, unlike the all-renter coverage rates in
Section~\ref{sec:twonyc}. The historical-school coverage in Chile uses pupils
linked to April registers. These targets and denominators are kept separate.
The published jobs differences average paired seed-block contrasts; intervals
resample the 20 blocks 10,000 times. The original GPT placement differences
have intervals $[-16.6,-14.4]$ and $[-18.9,-16.0]$ points. Anchor gains have
intervals $[11.9,13.9]$ and $[14.8,17.1]$ points. Repeated bootstrap draws
do not add markets or independent donor populations.

\paragraph{What coordination uses.}
Both anchored policies access the same entire feasible graph and capacities.
The assignment first maximizes cardinality and then reuse of displayed pairs.
An already displayed anchor changes nothing; otherwise one displayed posting
is replaced. GPT drops its last-ranked displayed pick; the random arm uses
its seeded drop rule. Neither construction uses a full-list partner label
to choose its anchor. Maximum feasible cardinality is an upper bound on DA
placement and supplies no individual retention guarantee.

\paragraph{What the anchor certifies.}
Write $\nu(G)$ for the maximum matching cardinality of a feasible slot graph
$G$. A displayed posting exposes all its feasible slots.
\begin{proposition}\label{prop:anchor-capacity}
Let $M^*$ be a maximum matching of $E$. Insert its assigned posting into each
recipient's menu if absent, replacing one displayed posting. If every
recipient of $M^*$ has a nonempty original menu, the resulting graph $D'$
satisfies $\nu(D')=\nu(E)$ and preserves each person's posting quota.
Nevertheless, a maximum-cardinality anchor that also maximizes reuse of
displayed pairs can raise displayed capacity while reducing DA placement.
This can occur with four applicants, unit seats, two options on every menu,
strict preferences and priorities, and only one changed menu.
\end{proposition}
\begin{proof}
Every pair of $M^*$ remains displayed or is inserted. Thus
$M^*\subseteq D'\subseteq E$, proving the capacity equality. Replacement of
one posting by a previously absent posting preserves each nonempty quota.

For the second claim, use the following preferences, written best first.
Options absent from a preference list are infeasible. Menus are sets; their
printed order is not an application order.
\begin{center}\small
\begin{tabular}{@{}llll@{}}\toprule
Applicant & Feasible preferences & Original menu & Edited menu\\\midrule
$a$ & $3,2,0,1$ & $\{0,1\}$ & $\{0,1\}$\\
$b$ & $2,3,0,1$ & $\{0,1\}$ & $\{0,1\}$\\
$c$ & $3,0,2$   & $\{0,3\}$ & $\{0,2\}$\\
$d$ & $0,1,3$   & $\{1,3\}$ & $\{1,3\}$\\\bottomrule
\end{tabular}\end{center}
Each option has one seat. Its priority order, highest first, is
\[
\begin{aligned}
0 &: c\succ a\succ b\succ d, & 1 &: d\succ a\succ b,\\
2 &: a\succ b\succ c,       & 3 &: c\succ d\succ b\succ a.
\end{aligned}
\]
The original display has capacity three because nobody sees option 2.
The anchor $(0,1,2,3)$, in applicant order, seats everyone and reuses three
displayed pairs; reusing four is impossible. It is therefore a valid
maximum-reuse maximum matching. Insert option 2 for $c$ and drop option 3.
This drop can be the recommender's last pick, even though $c$ prefers 3 in
the evaluation preferences. Only $c$'s menu changes. The original DA
assignment is $(0,\bot,3,1)$. After the edit, $c$ takes 0, displacing $a$;
neither $a$ nor $b$ can displace $d$ at 1. The new assignment is
$(\bot,\bot,0,1)$. Hence capacity rises from three to four while placement
falls from three to two, with one incumbent exit and no entrant.
\end{proof}
The example fixes a permitted anchor and drop rule. It does not assert failure
under every matching tie or drop convention, or estimate how frequently such
losses occur. The repository's coordinated-policy checks exercise this
instance against the archived anchor and DA implementations.

\paragraph{An edit-budget bound on individual exits.}
\begin{proposition}\label{prop:menu-edit-budget}
Fix strict applicant preferences, strict receiver priorities, feasibility and
responsive capacities. Compare applicant-proposing DA before and after edits
to $C$ applicants' menus. Let $X$ count recipients who become unmatched,
$Y$ count previously unmatched applicants who receive a seat, and
$\Delta=Y-X$ be the placement change. Then
\[
X\le C,\qquad Y\le C,\qquad X\le C-\max(\Delta,0).
\]
\end{proposition}
\begin{proof}
For a single edited applicant $i$, first run DA without $i$, obtaining
recipient set $R$. DA's rejection argument and applicant optimality imply
that its outcome is independent of the order of eligible proposals
\citep[Theorem~2]{gale1962}. We may therefore postpone $i$'s first proposal
until this smaller run terminates, then continue DA. Its previously unmatched
applicants have exhausted their menus. Along the ensuing rejection chain
there is at most one applicant still making proposals; a receiver replacing
a held applicant does not lose an occupied seat. The final recipient set is
therefore $R$, $R\cup\{i\}$, or
$(R\setminus\{r\})\cup\{i\}$ for some $r\in R$.
Apply this description separately to $i$'s old and new menus: their recipient
sets differ by at most one applicant in each direction. For $C$ edits,
change the menus sequentially. Every final exit or entry occurs at an
intermediate step, giving $X,Y\le C$. Finally $Y=X+\Delta$ gives the last
inequality. Responsive capacities admit the same rejection-chain argument.
\end{proof}
The bound concerns any-seat receipt, not posting or rank changes. Exits need
not belong to the edited group. The bound sums across disjoint markets but
does not cover endogenous priorities or changes to others' preferences.

In the frozen jobs panel, $N=5{,}760$, placement changes from $3{,}595$ to
$4{,}341$, and $C=1{,}205$. Thus $0\le X\le 459$; in replication,
$N=3{,}840$, placement changes from $2{,}437$ to $3{,}049$ and $C=920$,
giving $0\le X\le308$. These upper bounds are 7.969\% and 8.021\% of
all seeker--market observations. Counts are reconstructed from the released
means because every primary market has 24 seekers and equal weight.
These are outer intervals, not confidence intervals, observed exits, or
sharp bounds under this anchor. Actual transitions require individual records.

\paragraph{Available evidence and remaining computation.}
The aggregate source tables give anchored GPT placement rates of 75.3646\%
and 79.4010\%, compared with 78.1771\% and 81.4063\% for anchored random.
Subtracting these released means gives $-2.8125$ and $-2.00521$ points.
Their paired uncertainty cannot be recovered by subtracting marginal
intervals. The pending replay must evaluate random, GPT, anchored random
and anchored GPT on identical markets, retain assignment identities, and
count exits, entrants and rank changes relative to each specified baseline.
It must also measure post-anchor coverage, preserve each menu budget, verify
feasibility and complete cached responses, and report outcomes by wage group.
Apart from the outer bounds above, this manuscript contains no new result for
those endpoints. The current
integration has aggregate tables but lacks the licensed extract and private
reply caches required for that replay.

\paragraph{Timing and reproduction.}
The market design was frozen on November 15, 2026. Coverage and capacity/anchor
analyses followed the frozen-window outcomes but were specified before the
replication requests. The new coverage-after-anchoring, individual-access
and coordinated-baseline uncertainty analyses are follow-ups on both windows.
The original configurations and implementation are in
\texttt{configs/jobs\_frozen.json}, \texttt{configs/jobs\_replication.json}
and \texttt{experiments\_jobs/} at the repository root.
\texttt{python3 jobs/rebuild.py} verifies the four public source files against
\texttt{jobs/provenance.json}, then rebuilds Table~\ref{tab:jobs-coordination}
and \texttt{jobs/access-bounds.csv}. It runs no new markets or model calls;
postings, identifiers and reply caches are excluded.
